% Capítulo 14: Problemas Adicionales de Cálculo Vectorial

\subsection*{Ejercicios 1--3: Integrales impropias en varias variables, desarrollos en series infinitas de potencias y sumas de Euler}
\begin{enumerate}
    \subimport{ejercicio_01/}{ejercicio.tex}
    \subimport{ejercicio_02/}{ejercicio.tex}
    \subimport{ejercicio_03/}{ejercicio.tex}
\end{enumerate}

\subsection*{Ejercicios 4--7: Optimización geométrica de contornos, integrales sobre parte entera (función piso), operadores max y promedios integrales}
\begin{enumerate}
    \setcounter{enumi}{3}
    \subimport{ejercicio_04/}{ejercicio.tex}
    \subimport{ejercicio_05/}{ejercicio.tex}
    \subimport{ejercicio_06/}{ejercicio.tex}
    \subimport{ejercicio_07/}{ejercicio.tex}
\end{enumerate}

\subsection*{Ejercicios 8--10: Planos tangentes extremos, simplificaciones integrales iteradas tridimensionales e integrales impropias de arcos tangentes}
\begin{enumerate}
    \setcounter{enumi}{7}
    \subimport{ejercicio_08/}{ejercicio.tex}
    \subimport{ejercicio_09/}{ejercicio.tex}
    \subimport{ejercicio_10/}{ejercicio.tex}
\end{enumerate}

\subsection*{Ejercicios 11--12: Regiones sólidas definidas por proyecciones vectoriales y demostración del Ángulo Sólido mediante el Teorema de la Divergencia}
\begin{enumerate}
    \setcounter{enumi}{10}
    \subimport{ejercicio_11/}{ejercicio.tex}
    \subimport{ejercicio_12/}{ejercicio.tex}
\end{enumerate}
